\section{Hypergraph Category Axioms}\label{sec:axioms}

In a hypergraph category every object $X$ carries a special commutative
Frobenius algebra (SCFA, also called a Frobenius monoid) consisting of four morphisms
\begin{gather*}
  \mu_X \colon X \otimes X \to X,\quad
  \eta_X \colon I \to X,\\
  \delta_X \colon X \to X \otimes X,\quad
  \epsilon_X \colon X \to I,
\end{gather*}
subject to the axioms below, together with coherence conditions that fix the
SCFA on every tensor product $A \otimes B$ in terms of those on $A$ and $B$.
Throughout, $\sigma_{X,Y}\colon X\otimes Y\to Y\otimes X$ is the symmetry of the underlying
symmetric monoidal category, the wire crossing, and $;$ denotes composition in diagrammatic order.

\paragraph{Monoid.}
\begin{align}
  (\mu_X \otimes \mathrm{id}_X) ; \mu_X
    &= (\mathrm{id}_X \otimes \mu_X) ; \mu_X
    \label{eq:ax-assoc}
    \tag{assoc}\\
  (\eta_X \otimes \mathrm{id}_X) ; \mu_X
    &= \mathrm{id}_X
    = (\mathrm{id}_X \otimes \eta_X) ; \mu_X
    \label{eq:ax-unit}
    \tag{unit}
\end{align}

\paragraph{Comonoid.}
\begin{align}
  \delta_X ; (\delta_X \otimes \mathrm{id}_X)
    &= \delta_X ; (\mathrm{id}_X \otimes \delta_X)
    \label{eq:ax-coassoc}
    \tag{coassoc}\\
  \delta_X ; (\epsilon_X \otimes \mathrm{id}_X)
    &= \mathrm{id}_X
    = \delta_X ; (\mathrm{id}_X \otimes \epsilon_X)
    \label{eq:ax-counit}
    \tag{counit}
\end{align}

\paragraph{Frobenius.}
\begin{align}
  (\mathrm{id}_X \otimes \delta_X) ; (\mu_X \otimes \mathrm{id}_X)
    &= \mu_X ; \delta_X\notag\\
    &= (\delta_X \otimes \mathrm{id}_X) ; (\mathrm{id}_X \otimes \mu_X)
    \label{eq:ax-frob}
    \tag{Frob}
\end{align}

\paragraph{Commutativity and cocommutativity.}
\begin{align}
  \sigma_{X,X} ; \mu_X &= \mu_X
    \label{eq:ax-comm}
    \tag{comm}\\
  \delta_X ; \sigma_{X,X} &= \delta_X
    \label{eq:ax-cocomm}
    \tag{cocomm}
\end{align}

\paragraph{Speciality.}
\begin{align}
  \delta_X ; \mu_X &= \mathrm{id}_X
    \label{eq:ax-special}
    \tag{special}
\end{align}

\paragraph{Monoidal coherence.}
The SCFA on $A \otimes B$ is determined by those on $A$ and $B$:
\begin{align}
  \mu_{A \otimes B}
    &= (\mathrm{id}_A \otimes \sigma_{B,A} \otimes \mathrm{id}_B)
       ; (\mu_A \otimes \mu_B)
    \label{eq:ax-coh-mu}
    \tag{coh-$\mu$}\\
  \eta_{A \otimes B}
    &= \eta_A \otimes \eta_B
    \label{eq:ax-coh-eta}
    \tag{coh-$\eta$}\\
  \delta_{A \otimes B}
    &= (\delta_A \otimes \delta_B)
       ; (\mathrm{id}_A \otimes \sigma_{A,B} \otimes \mathrm{id}_B)
    \label{eq:ax-coh-delta}
    \tag{coh-$\delta$}\\
  \epsilon_{A \otimes B}
    &= \epsilon_A \otimes \epsilon_B
    \label{eq:ax-coh-eps}
    \tag{coh-$\epsilon$}
\end{align}
The monoidal unit $I$ carries the trivial SCFA, with $\mu_I = \eta_I = \delta_I = \epsilon_I = \mathrm{id}_I$.
